\section{The separated-link coefficient on actual compact Weil tests}
\label{sec:xi4-actual-Weil}

The full supporting calculation in
Section~\ref{the-first-separated-link-connected-electric-covariance}
uses the original finite-box Hamiltonian convention of
Kogut--Susskind~\cite{kogutSusskind1975}. Its fourth coefficient is
derived there, not attributed to that historical paper. Here we carry
that coefficient into the \emph{classical} Weil form with its explicit
inverse, and then evaluate its actual test functions. The arithmetic
form is still \eqref{eq:rh-full-weil-test}.

\subsection{The retained physical parameters and arithmetic observable}

For integer $L\ge2$ retain the $N=3(2L)(2L+1)^2$ physical links and
$M=3(2L)^2(2L+1)$ elementary faces of the original open box. Write
\[
 \begin{gathered}
 H=\kappa(H_0-\xi S)+2b_{\rm mag}MI,\qquad
 H_0=\sum_e E_e,\quad S=\sum_p W_p,\\
 \kappa=2g^2/a_{\rm lat},\qquad
 b_{\rm mag}=1/(2g^2a_{\rm lat}),\quad \xi=1/(4g^4).
 \end{gathered}
\]
The subscripts retain the lattice spacing and magnetic coupling,
whose source names are $a,b$; they do not identify them with the
polynomial target coordinates. More precisely, on the product of the
positive parameter plane with the original polynomial coordinate
space, the map is
$(a_{\rm lat},g,q)\mapsto(a_{\rm lat},g,F(q))$.
It has derivative $\operatorname{diag}(I_2,DF)$ and arithmetic
projection $F_1(q)/2$. Both physical parameters are unchanged.
The positive parameter chart has the explicit inverse
\[
 (a_{\rm lat},g)\mapsto(\kappa,\xi),\qquad
 (\kappa,\xi)\mapsto
       (1/(\kappa\sqrt{\xi}),(4\xi)^{-1/4}),\qquad b_{\rm mag}=\kappa\xi.
\]
Substitution gives both identities. The original generators
$T_a=-i\sigma_a/2$ satisfy
$-\operatorname{tr}(T_aT_b)/2=\delta_{ab}/4$.
Thus the original metric's orthonormal frame is $2T_a$ and
$-\Delta^c=4E_e$; its electric coefficient is $\kappa/4$.

Let $\psi_\xi$ be the actual positive unit vacuum, with energy
$2b_{\rm mag}M+\kappa e(\xi)$, and set
$A_\xi=H_0-\xi S-e(\xi)I$ and
$v_e=(E_e-\langle\psi_\xi,E_e\psi_\xi\rangle)\psi_\xi$.
For distinct links sharing no face, the supporting proof gives
\begin{equation}
 \lim_{\xi\to0}\xi^{-4}
       \langle v_e,h(A_\xi)v_f\rangle=a_{ef}\mathcal D(h),\qquad
 \mathcal D=c_0\delta_3+\delta'_3/336+
                         c_1\delta_{9/2}+c_2\delta_{13/2},
 \label{eq:xw-source}
\end{equation}
for bounded continuous $h$ differentiable at $3$, where
\[
 c_0=-59/275184,\qquad c_1=1/1296,\qquad c_2=3/132496.
\]
The integer $a_{ef}$ is the exact ordered adjacent-face count there.
In particular $\delta'_3(h)=-h'(3)$,
$\mathcal D(1)=127/219024$, and $\mathcal D(\lambda)=0$.
The zero first moment retains the band derivative and both higher
spin channels; it does not remove the positive zeroth moment.
For this general class of $h$, \eqref{eq:xw-source} gives a
little-$o(\xi^4)$ remainder. An even $O(\xi^6)$ remainder requires
the further analytic dependence proved for the specified source
observables, not just bounded continuity of $h$.

Here is the direct connection to the unchanged arithmetic function.
Put $k(v)=\Phi(|v|/2)/2$, with $\Phi$ from
\eqref{eq:cq-realheat}. Changing variables in the cosine integral gives
\[
 Z_t(s)=4\int_{\mathbb R}k(v)e^{tv^2/4+isv}\,dv.
\]
The superexponential estimate proved below
\eqref{eq:cq-realheat} absorbs every fixed quadratic, linear and
polynomial weight. Hence for every complex $t,s$ the function
$h(q)=Z_t(s+q)$ and all its real $q$ derivatives are bounded:
the factor $e^{iqv}$ has modulus one. For real coupling the spectral
theorem and Fubini, applied to each pair of Hilbert-space vectors,
give the actual bounded operator
\[
 Z_t(s+A_\xi)=4\int_{\mathbb R}
          k(v)e^{tv^2/4+isv}e^{ivA_\xi}\,dv.
\]
The integral is strong on each vector and has operator norm at most
$4\int|k(v)|e^{\operatorname{Re}(t)v^2/4+|\operatorname{Im}s||v|}dv$.
This dominating integral also proves norm-entire dependence on
$t,s$ by integrating their absolutely convergent Taylor series.
Applying \eqref{eq:xw-source} therefore yields the actual coefficient
$a_{ef}\mathcal T_{\mathcal D}Z_t$, where
\begin{equation}
 (\mathcal T_{\mathcal D}f)(s)
 =c_0f(s+3)-f'(s+3)/336+
                          c_1f(s+9/2)+c_2f(s+13/2).
 \label{eq:xw-TD}
\end{equation}
In particular its initial value contains the actual $\Xi$ and its
derivative at these translated arguments. No freely chosen entire
function has replaced the Riemann datum.

\subsection{Exact compact-support map, inverse and all jets}

Use the original minus-sign Fourier transform
$\mathscr Fg(s)=\int g(u)e^{-isu}du$. Define
\begin{align}
 m(v)&=e^{3iv}(c_0-iv/336)+c_1e^{9iv/2}+c_2e^{13iv/2},\notag\\
 w(u)&=m(-u)=e^{-3iu}(c_0+iu/336)
                             +c_1e^{-9iu/2}+c_2e^{-13iu/2}.
 \label{eq:xw-multiplier}
\end{align}
Translation and differentiation under the integral prove the exact
typed correspondence
\begin{equation}
 \mathcal T_{\mathcal D}\mathscr F
       =\mathscr F M_w,\qquad M_wg(u)=w(u)g(u).
 \label{eq:xw-transform}
\end{equation}
The sign reversal in $w=m(-u)$ is essential.

For all real $v$ there is a uniform rational lower bound
$|m(v)|\ge\delta=575/1752192$. Here is its full elementary proof.
For $|v|\le1/2$, factoring $e^{3iv}$ and using
$\cos x\ge1-x^2/2$ gives
\[
 \Re(e^{-3iv}m(v))\ge
 c_0+(23/32)c_1-(17/32)c_2=575/1752192.
\]
The cosine inequality follows by twice integrating
$1-\cos x\ge0$. For $|v|\ge1/2$, the reverse triangle inequality gives
\[
 |m(v)|\ge |v|/336-(|c_0|+c_1+c_2)
       \ge10291/21464352>575/1752192.
\]
These ranges cover the whole line, proving the bound without
removing a Fourier point or choosing a branch.

For every compact interval $I$ and $r=2$ or $\infty$, $M_w$ is
an automorphism of the functions in $C_c^r(\mathbb R)$ supported in
$I$, with inverse $M_{1/w}$ and \emph{exactly} unchanged support.
Indeed multiplication by a nowhere-zero smooth function preserves
the nonzero set. Leibniz's rule bounds each derivative seminorm
of $wg$ by finitely many of those of $g$. The inverse has the same
property: differentiate $w(1/w)=1$ recursively and use $|w|\ge\delta$
on $I$. Both pointwise compositions are the identity. Thus the
map also gives a continuous automorphism of $C_c^\infty$ with its
usual fixed-support inductive-limit topology, and of its Fourier
image with the transported topology.

The full inverse jet coordinates, rather than an identification of
pointwise zero sets, are particularly useful. For an output
$G=\mathscr Fg$ define
\begin{equation}
 \mathcal L_{z,n}(G)_r
     =\int_I(-iu)^r\,\frac{g(u)}{w(u)}e^{-izu}du,
       \quad 0\le r<n.
 \label{eq:xw-inverse-jets}
\end{equation}
Then
$\mathcal L_{z,n}\mathcal T_{\mathcal D}\mathscr Ff
 =((\mathscr Ff)^{(r)}(z))_{r<n}$.
This follows by substituting $g=wf$ and differentiating an integral
over a fixed compact set. Its absolute value is bounded by
$\delta^{-1}\int_I |u|^r|g(u)|e^{|\operatorname{Im}z||u|}du$,
which proves continuity.
Conversely the forward coordinates are, exactly,
\[
 (\mathcal T_{\mathcal D}\mathscr Ff)^{(r)}(z)
  =c_0\widehat f^{(r)}(z+3)-\widehat f^{(r+1)}(z+3)/336
    +c_1\widehat f^{(r)}(z+9/2)+c_2\widehat f^{(r)}(z+13/2).
\]
Thus a zero and all its multiplicity data are transported by
\eqref{eq:xw-inverse-jets}; no equality of untransported evaluations
is imposed.

For completeness, each finite ordinary jet map onto $\mathbb C^n$
is surjective on the compact smooth Fourier image. Choose a
nonnegative nonzero smooth bump $\chi$ inside an interval with
nonempty interior. The $n$ polynomials times
$e^{izu}\chi(u)$ give a moment matrix
$\bigl(\int u^{j+k}\chi(u)du\bigr)_{j,k<n}$, times invertible
diagonal powers of $-i$. The moment matrix is positive definite:
its quadratic form is $\int|p(u)|^2\chi(u)du$, which vanishes only
when the polynomial vanishes on an interval and hence identically.
Matrix inversion gives every prescribed jet supported in
$\operatorname{supp}\chi$. To retain exactly this support even
for the zero jet, use the same moment matrix to construct the
monic degree-$n$ polynomial $p_n$ orthogonal to degrees below $n$
for weight $\chi$. Adding $p_n$ to the degree-less-than-$n$
interpolant changes none of its prescribed moments. The resulting
polynomial is nonzero and has only finitely many real roots, so
its product with $e^{izu}\chi$ has exactly $\operatorname{supp}\chi$.
Composition with $M_w$ proves surjectivity of
\eqref{eq:xw-inverse-jets}. If $\mathcal J$ is the ordinary jet
map, then $\mathcal T_{\mathcal D}\ker\mathcal J=\ker\mathcal L$;
the induced quotient isomorphism and its inverse are the displayed
maps on representatives. It is the identity in their $\mathbb C^n$
jet coordinates. The entire recovered function also retains every
analytic-unit coefficient: at a zero of order $n$ they are
$\mathcal L_{z,n+j+1}(G)_{n+j}/(n+j)!$, for $j\ge0$.

The endpoint constraints in Section~\ref{sec:endpoint-Weil} therefore
transport to coordinates $1$ through $4$ of $\mathcal L_{\sigma,5}$,
leaving coordinate $0$ unrestricted. This describes their actual image
under $M_w$, not another four constraints chosen on the output.
The heat multiplier $e^{tu^2/4}$ commutes with $w$ and $1/w$ on
every compact-support space. Conjugating by $\mathscr F$ gives
commutation with the full complex-time heat group of
\eqref{eq:cq-testheat}, in both directions.

\subsection{The full arithmetic pullback and its evaluated channels}

For arbitrary compact tests $f,g$ the transformed correlation is
\begin{equation}
 H^w_{f,g}(x)=\int_{\mathbb R}
       w(v)g(v)\,\overline{w(v-x)f(v-x)}\,dv.
 \label{eq:xw-correlation}
\end{equation}
Substitute this expression, its Fourier transform and its values at
every $\pm\log n$ into \eqref{eq:rh-full-weil-test}. This is the
complete pulled-back form $W_w(f,g)=W(M_wf,M_wg)$.
It retains both poles, the whole gamma integral and all prime powers.
The map on ordered test pairs is the bijection
$(f,g)\mapsto(wf,wg)$ with inverse $(f/w,g/w)$; hence it gives the
exact correspondence between strict negative witnesses of $W_w$
and those of $W$. It is not a claim that $W_w(f,g)=W(f,g)$.
In particular the distributional coefficient $\mathcal D$ itself
is not substituted for this arithmetic pairing.

At the previously retained endpoint $b_{\rm pol}=8$, $\sigma=2$,
start with $f_0=e^{2iu}b_4(u)$. The source image is exactly
\[
 wf_0=c_0f_1+f_2/336+c_1f_3+c_2f_4,\quad
 (f_1,f_2,f_3,f_4)
  =(e^{-iu}b_4,\;iu e^{-iu}b_4,\;
                     e^{-5iu/2}b_4,\;e^{-9iu/2}b_4).
\]
Every test is supported in $[-2,2]$. We evaluate the entire
five-by-five Gram matrix, not just the displayed linear combination.
The functions are independent: on an open interval where $b_4>0$,
a vanishing combination is an exponential polynomial with distinct
frequencies $2,-1,-5/2,-9/2$ and a degree-one coefficient at $-1$.
Successive application of the differential factors
$\partial_u-i\sigma_j$ annihilating the other terms isolates each
highest polynomial coefficient; descending its degree proves
all coefficients zero.

Here the gamma integral was independently evaluated in frequency
coordinates. Set $\varphi(z)=\operatorname{sinc}(z/2)^4$.
The five Fourier functions are
\[
 (\Psi_0,\ldots,\Psi_4)(z)
   =(\varphi(z-2),\varphi(z+1),-\varphi'(z+1),
                        \varphi(z+5/2),\varphi(z+9/2)).
\]
They are real on the real axis. Thus the two pole terms are
$2\Re(\Psi_i(i/2)\Psi_j(i/2))$, the gamma integrand is
$\Psi_i(t)\Psi_j(t)$ times the exact weight in
\eqref{eq:rh-full-weil-test}, and the prime terms use the real
parts of $f_j*f_i(\log n)$. This last convolution identity follows
because each $f_i^\star=f_i$. Its four cubic spline pieces and their
shifted knots are integrated without approximation of the phase.
All prime powers $n<e^4$ are included.

Here is the explicit bound for the discarded \emph{integration
range}, which is added back as an interval error.
Take $T=96$, $A=9/2$, $\ell=T-A$.
For $|t|\ge T$ one has $|\Psi_i(t)|\le C_i/(|t|-A)^4$,
where $C_i=16$ except $C_2=32+64/\ell$.
Indeed $\varphi(q)=16\sin^4(q/2)/q^4$ and
$|\varphi'(q)|\le32/|q|^4+64/|q|^5$ on the real line.
The digamma bound already proved in
\eqref{eq:ew-continuity}, keeping its logarithmic bound, gives
\[
 |\text{gamma weight}(t)|
 \le\{\gamma_E+8+\log\pi+\log(1+2|t|)\}/(2\pi).
\]
Concavity bounds the last logarithm, for $t\ge T$, by
$\log(1+2T)+(t-T)/(T+1/2)$. Direct integration on the two tails
therefore bounds the absolute omitted matrix entry by
\begin{equation}
 \frac{C_iC_j}{\pi}
 \left\{\frac{\gamma_E+8+\log\pi+\log(1+2T)}{7\ell^7}
                +\frac1{42(T+1/2)\ell^6}\right\}.
 \label{eq:xw-gamma-tail}
\end{equation}
This is an error enclosure, not deletion of a tail or truncation
of a zero sum.

The saved 192-bit Arb calculation, with $2^{-100}$ quadrature
tolerances, proves all five leading principal determinants positive.
All cross entries are retained. Its source vector and optimal
subtraction of the seed obey the deliberately outward bounds
\begin{align*}
 6.63\,10^{-14}&<W(wf_0,wf_0)<6.77\,10^{-14},\\
 4.70\,10^{-14}
 &<W(wf_0,wf_0)-|W(f_0,wf_0)|^2/W(f_0,f_0)
                                                    <4.90\,10^{-14}.
\end{align*}
The second expression is the exact minimum over multiples of $f_0$,
by completing the square, with minimizer
$wf_0-f_0W(f_0,wf_0)/W(f_0,f_0)$.
The certificate also verifies overlap of its $(0,0)$ entry with
the earlier independently evaluated real-space gamma formula.
The full entries, quadrature pieces, prime correlations and error
bounds are retained in \texttt{xi4\_weil\_results.json};
\texttt{verify\_xi4\_weil.py} reproduces them under a one-worker
5,000,000,000-byte process-tree ceiling. No zero ordinates enter.
The read-only \texttt{replay\_xi4\_weil.py} independently computes
all $360$ prime correlations by exponential-polynomial
antiderivatives instead of quadrature. On each exact spline cell,
for a polynomial $P$ of degree $d$ and $a\ne0$, it uses
\[
 \int P(v)e^{av}dv
   =e^{av}\sum_{r=0}^{d}(-1)^rP^{(r)}(v)/a^{r+1}.
\]
Differentiation telescopes to the integrand; at $a=0$ it integrates
each polynomial monomial directly. All $15$ independently assembled
matrix entries overlap the original enclosures. Independent
interval LDL elimination proves all five pivots positive both for
the saved matrix and for this reconstructed matrix. The replay
also recomputes the pole terms, the sum of the saved $192$ gamma
integration cells, the explicit tail bound, and both displayed
quadratic values. The frequency-side gamma integrations remain
trusted Arb integrals: this replay does not mislabel their saved
values as a second integration. Serialization enlarges the
replayed Schur enclosure to $[4.8\,10^{-14}\ \pm4.02\,10^{-16}]$;
the displayed bounds contain it as well as the original enclosure.

These tests are $C_c^2$ splines. Write $S_\varepsilon$ for the
moment-preserving smoothing at $\sigma=2$ in
Section~\ref{sec:endpoint-Weil}. For each computed output test $h_i$,
the exact smooth approximation retaining its inverse jet
coordinates is
\[
 h_{i,\varepsilon}=M_w S_\varepsilon M_{1/w}h_i .
\]
The reciprocal map first gives a compact $C^2$ input.
$S_\varepsilon$ preserves its moments of orders zero through four
at $\sigma$, and multiplication back by $w$ gives the displayed
transported coordinates unchanged. On one common enlarged compact
interval both multiplier maps are bounded in $C^2$ by the
product-rule estimates above. Thus $h_{i,\varepsilon}\to h_i$
in $C^2$. Applying the complete Weil continuity bound to every
pair proves convergence of the entire Gram matrix, so each
strictly positive leading determinant stays positive for all
sufficiently small $\varepsilon$. Directly smoothing $h_i$ and
then applying $M_w$ would instead converge to $M_wh_i$; that is
not the family calculated here. The positive five-dimensional result
is not positivity of the infinite compact-test space. It supplies
no RH counterexample, but evaluates the actual source multiplier
and its full arithmetic cross terms.

\input{proofs/C_xi4_original_source}
